\section{Divergence}
\label{app:divergence}

\paragraph{Methods.}
The divergence study scores each prompt at four readouts (3-opt, 2-opt, greedy-,
and sampled-thinking) under the per-condition gold. The sampled-thinking variant
draws $N$ generations $g_i \sim \LLM(x)$, parses each to $r_i$, and over the
$M \le N$ parsed draws forms
\begin{equation}
  p^{\mathrm{th}}_y \;=\; \frac{1}{M}\sum_{i=1}^{M} \mathbf{1}[\, r_i = y \,],
  \label{eq:pth}
\end{equation}
with prediction $\hat{y}^{\mathrm{th}} = \argmax_y p^{\mathrm{th}}_y$ (ties resolve
to $u$). To characterize the reasoning distribution we use Shannon entropy and
divergence from correct abstention,
\begin{equation}
  H(p) = -\sum_{y} p_y \log p_y,
  \qquad
  \mathrm{JS}\!\left(p^{\mathrm{th}} \,\|\, \delta_u\right),
  \label{eq:entropy-js}
\end{equation}
where $\mathrm{JS}$ is the Jensen--Shannon divergence (low when free-form reasoning
concentrates on $u$).

\paragraph{Results.}
\begin{itemize}[leftmargin=*,topsep=2pt,itemsep=1pt]
  \item Scored set: $43$ items of $47$ prompts ($4$ unparseable thinking draws);
    $14$ ambiguous $+\,29$ disambiguated.
  \item Ambiguous accuracy by readout: sampled thinking $86\%$ ($12/14$);
    three-option $79\%$ ($11/14$); greedy-thinking $64\%$ ($9/14$); two-option
    undefined (no unknown option).
  \item Disambiguated accuracy by readout: two-option $59\%$ ($17/29$); sampled
    thinking $48\%$ ($14/29$); greedy-thinking $41\%$ ($12/29$); three-option
    $28\%$ ($8/29$).
  \item Ambiguous default thinking mix: $84\%$ unknown, $9\%$ target, $7\%$ other;
    $11$ of $14$ items sit at $[0,0,1]$ (JS-divergence $0$).
  \item Mean JS-divergence from $\delta_u$: $0.092$; three items diverge (one onto
    \emph{other} at $\mathrm{JS}=\ln 2 = 0.693$, two toward target at $0.216$ and
    $0.380$).
\end{itemize}

\begin{figure}[htbp]\centering
  \begin{subfigure}{0.49\linewidth}\centering
    \plotsub{../out/sesgo/divergence/Qwen3-0.6B/plots/accuracy_by_readout.png}{accuracy_by_readout.png}{\linewidth}
    \caption{Per-condition accuracy across the four readouts.}
    \label{fig:divergence-acc}
  \end{subfigure}\hfill
  \begin{subfigure}{0.49\linewidth}\centering
    \plotsub{../out/sesgo/divergence/Qwen3-0.6B/plots/role_mix.png}{role_mix.png}{\linewidth}
    \caption{Ambiguous default role mix: $84\%$ unknown, $9\%$ target, $7\%$
      other.}
    \label{fig:role-mix}
  \end{subfigure}
  \caption{\textbf{Divergence} on \texttt{Qwen3-0.6B} ($43$ scored items of $47$
    prompts---$4$ had unparseable thinking draws---$14$ ambiguous $+\,29$
    disambiguated): per-condition accuracy by readout (left) and the ambiguous
    default role mix (right).}
  \label{fig:role-geo}
\end{figure}

\plotfig{../out/sesgo/divergence/Qwen3-0.6B/plots/thinking_belief_agreement_scatter.png}{thinking_belief_agreement_scatter.png}{0.92\linewidth}{\textbf{Does thinking change the answer?} Per ambiguous item, the chance the model abstains (`unknown') \emph{before} reasoning (a direct read) vs.\ \emph{after} reasoning step by step (averaged over its free-form tries); a point on the diagonal means thinking left the belief unchanged. Pre- and post-thinking abstain beliefs only loosely agree ($r{=}0.50$; just $44\%$ of items sit within $0.15$ of the diagonal), and on balance thinking pulls the model \emph{away} from the safe abstention (mean post$-$pre $={-}0.11$).}{fig:thinking-agreement}

\paragraph{Dynamics: forking-paths.}
Following the forking-paths method, we fork the greedy reasoning path of one
representative ambiguous item at every token position ($60$ positions, $40$ forks
each, a $60$-sample prior), parse each continuation to an outcome, and track the
outcome distribution $O_t$ along the path; a Bayesian change-point posterior
$p(\tau{=}t\mid y)$ tests for a single commitment token, and a branching tree
summarizes the root $\to$ trunk $\to$ identity-branch outcome mix.
\begin{itemize}[leftmargin=*,topsep=2pt,itemsep=1pt]
  \item $O_t$ drifts along the path (a stereotyped-group excursion near
    $t{\approx}40$, then a swing toward abstention) but the change-point test finds
    \emph{no} significant single commitment token (Bayes factor $0.023$,
    $p(m{=}0){=}0.98$): on this small model the answer accretes gradually rather
    than snapping at one token---an honest negative.
  \item The strongest fork is at $t{=}42$ (re-sampling there most diverts the
    outcome): pull $0.284$, drift $0.341$, potential $0.285$.
\end{itemize}

\begin{figure}[htbp]\centering
  \begin{subfigure}{0.62\linewidth}\centering
    \plotsub{../out/sesgo/forking/Qwen3-0.6B/plots/forking_commit_dynamics.png}{forking_commit_dynamics.png}{\linewidth}
    \caption{Outcome distribution $O_t$ over reasoning tokens (top) and the
      change-point posterior (bottom; no single deciding token).}
    \label{fig:forking-dynamics}
  \end{subfigure}\hfill
  \begin{subfigure}{0.36\linewidth}\centering
    \plotsub{../out/sesgo/divergence/Qwen3-0.6B/plots/branching_tree.png}{branching_tree.png}{\linewidth}
    \caption{Branching tree: root $\to$ trunk $\to$ named / abstain branches.}
    \label{fig:branching-tree}
  \end{subfigure}
  \caption{\textbf{Dynamics (forking-paths)} on \texttt{Qwen3-0.6B}: how the answer
    forms across the reasoning path (left) and the branching outcome tree (right).}
  \label{fig:dynamics}
\end{figure}
